\chapter{Introduction\label{cha:chapter1}}

Warehouses do not audit themselves, and the fleet built to help them do so in this thesis coordinates without any single robot, or any external process, ever holding authority over the others. This chapter discusses why that constraint condition matters before this thesis spends eight more chapters building out a system around it. Section~\ref{sec:ch1_motivation} recounts the motivation for casting inventory auditing as a robotics problem in the first place, with an emphasis on the coverage-versus-disruption trade structure underlying manual audits. Section~\ref{sec:ch1_tension} narrows this motivation to the particular tension this thesis seeks to address: coordination in a decentralized manner offers resilience to single points of failure, but the empirical record of that trade is less rosy than the folklore suggests. Section~\ref{sec:ch1_scope} takes that question back to the domain of a specific fleet of three robots, identifying a bounded claim about them this thesis aims to demonstrate, and Section~\ref{sec:ch1_contributions} lists the seven components used to do so. Finally, Section~\ref{sec:ch1_organization} closes the chapter by briefly describing each of the eight chapters that make up the remainder of this work.

\section{Motivation: Manual Inventory Auditing and the Case for Autonomous Warehouse Robotics\label{sec:ch1_motivation}}

At its base level, a warehouse audit is a force balance - a physical shelf contains some amount of inventory, while the corresponding digital record contains some other amount, and the disconnect between the two, if not reconciled, cascades down the supply chain in the form of misplaced stock, improper safety stocks, stockouts, and rerouted orders. The most direct way to close the gap is to send someone down every aisle with a scanning device, but that process scales poorly, not only in terms of time but also in terms of space, as auditors and pickers must share aisle access.

This tension between coverage and disruption underpins the frequency of audits, and larger facilities with higher racks amplify the problem, as verifying stock on a shelf six meters above the floor requires either a forklift, a ladder, or a willingness to climb, none of which are conducive to continuous operations. It is this tension that has made auditing a perennial candidate for mechanization and automation in recent years, and a decentralization of the audit process has appeared as a natural way to reduce the cost in operational space while retaining coverage, if not ensuring it~\cite{pallottino2026warehouses}.

The turn to alternatives to manual, single-vehicle data collection has seen the rise of both ground-based automated mobile robots participating in intralogistics transport and specialized inventory inspection platforms that take on the task of stock verification~\cite{bernardo2022survey}. These latter platforms have seen an increased focus on aerial support systems, as drones can inspect higher shelves without requiring additional equipment, and a synthesis of over a hundred works on drone-assisted logistics and inventory management highlights the overall lack of research considering multi-vehicle operations, where most works center around single-vehicle simulations in complex, dynamic warehouses, with few works dealing with actual implementations of fully autonomous multi-agent operations in realistic settings~\cite{pore2026drone}. Multi-agent coordination and safety are, therefore, a known issue, with most works pointing to a lack of relevant literature.

This thesis turns to a similar challenge but highlights a different set of constraints. By choosing a ground-based fleet of differential-drive robots, it replaces potential limitations in vertical reach with challenges in spatial coordination, asking whether such a system can manage the audit task in a warehouse environment. In this thesis, no audit is ever performed by a single robot, and no robot is ever given absolute authority over the process, as time spent waiting for permission to act is time spent blocking potential throughput. This coordination challenge is discussed in detail in Section~\ref{sec:ch1_tension} and is the reason why most works on multi-robot inventory auditing operate in warehouses without high-level shelving or other structures that penalize multi-agent coordination.

\section{The Coordination-Resilience Tension: Why Decentralization Complicates Fault Tolerance\label{sec:ch1_tension}}

A dispatcher provides coordination cheaply. One process keeps the map up-to-date, knows all the robots' locations, and forbids sending robots into the same aisle. The obvious appeal of this approach is that conflicts are solved by a single decision, easily auditable and understood by humans, and the process itself is straightforward to implement and maintain. It becomes considerably less attractive as soon as a single process gets killed in the middle of the update, loses connection to the robots, or simply gets overwhelmed trying to track the state of all agents in real-time. A fleet with a single point of coordination can afford to lose it, but a warehouse filled with shelf-storage does not provide the space needed to host an additional dispatcher robot as a backup.

The alternative, distributing the coordination logic between the robots, removing the single point of contention, has to negotiate its own set of challenges. The comparison of centralized and decentralized strategies, executed on actual hardware and evaluated against each other, not theoretical best-case scenarios, found the trade-off not as beneficial for some real-time applications as one could expect from the textbook models. The fully decentralized coverage algorithm demonstrated its ability to gracefully recover from failures and deal with a higher than average number of simultaneous faults in a scenario where robot failures were simulated, however, the increase in performance was relatively modest. It was enough to offset the difference in performance between centralized and decentralized approaches in most cases, except for a nearly theoretical level of failures, when the robot team executing a decentralized algorithm outperformed the one executing a centralized task allocation~\cite{jamshidpey2025centralization}. Decentralization, therefore, comes with a price of reduced maximum performance and a much more complex set of design constraints, especially in the domain of fault tolerance. The latter consideration takes its place in the wider discussion of task allocation, as a specific case of the broader challenges, which are, for the field of multi-robot systems, non-trivial and not yet settled.

The problems of decentralized fault-tolerant control of a robot formation are far from being solved; recent research into the subject is still trying to find common patterns for implementing fault detection and isolation, as an example, and is far from finding a solution applicable to most common actuating faults and sensors~\cite{elsayed2024decentralized}. Similar challenges are posed by the decentralized approaches to task allocation, spatial mutual exclusion, and power management, which, again, have their centralized counterparts offering better, well-defined solutions in some cases.

This thesis, therefore, operates within a limited domain of decentralized coordination algorithms for a small group of ground robots to demonstrate the possibility of a trade-off between the two approaches. It operates on actual hardware, where robots cooperate to navigate a warehouse with moderately narrow aisles. It limits the consideration to a homogeneous group of three robots, with the promise of no single point of coordination, even in the single point of contention failure case. It does not attempt to operate on the global scale of the multi-robot coordination problem, which is, to this day, an unsolved challenge of autonomous systems design and control. It asks a narrower question of whether a small fleet of robots, with peer-to-peer communication channel, without a fail-safe point of coordination, can continue performing its maintenance tasks, detect and respond to failures, and continue auditing the warehouse with minimal interruption to its operations, and at which cost, in terms of reduced maximum performance. The question is further formalized in Chapter~\ref{cha:chapter4}, answered from two sides in Chapters~\ref{cha:chapter5} and~\ref{cha:chapter6}.

\section{Research Objectives and Scope\label{sec:ch1_scope}}

The tension described in Section~\ref{sec:ch1_tension} is precisely the research objective of this thesis, which the following chapters are organized to address rather than merely describe. The objective conveniently decomposes into two subgoals, which correspond to the two operating regimes a fleet must be prepared to handle: operating under nominal conditions plus recovering gracefully from the loss of one of the agents during a mission. This thesis therefore asks whether a small fleet of ground robots, operating in a fully decentralized manner, is capable of (i) completing the warehouse audit task with throughput scaling approximately linearly in the fleet size under nominal conditions and (ii) reacting appropriately to the loss of one of the members, redistributing the load, and completing the audit to the end without the aid of any central coordination, not even temporary.

The system described in this thesis is intentionally minimalist but effective, as designating claims about decentralized coordination as true is challenging, and minimality helps reason about the system at large. Three TurtleBot3 Waffle differential-drive robots coordinate over ROS~2 in a peer-to-peer fashion, using topics rather than a dispatcher, and operate to audit 142 discrete points placed in twenty-two storage units in a mock warehouse. Coordination decisions are made exclusively by agents using the information provided to them, which is a critical feature even when considering the allocation decisions in Chapter~\ref{cha:chapter5} and routing decisions in Chapter~\ref{cha:chapter6} made after one of the robots fails.

There are two additional limitations imposed that deserve explicit mentioning. First, the autonomy stack used in the fleet, including the energy model, does not implement a return-to-dock procedure; the robots are effectively allowed to discharge fully before resuming operations. This choice is entirely defensible, as described in Chapter~\ref{cha:chapter4}, but it is also explicitly identified as a limitation in the broader context of warehouse operations in Chapter~\ref{cha:chapter9}. Second, and more importantly, all the results are obtained in a Gazebo Harmonic simulation running on a single WSL2 instance, rather than actual hardware. The discrepancy between the real and simulated robot dynamics is a known limitation this thesis does not address directly.

\section{Contributions of the Thesis\label{sec:ch1_contributions}}

Driven by the motivating conditions of Section~\ref{sec:ch1_scope}, this thesis introduces and evaluates the following contributions to a decentralized fleet of robots performing warehouse auditing tasks:

\begin{enumerate}
    \item \textbf{A deterministic, cardinal-direction sweep allocator with same-side-lock and same-item priority rules.} By forgoing the opportunity to plan task allocations as an online optimization, the fleet is able to encode the next best target as a decision tree based on region ownership, approach cardinality, and side-completion status. The loss in optimization power is mitigated by a decision procedure which is both nearly instantaneous to evaluate and provably consistent across robot instances in the same local state---a property suggested by Chapter~\ref{cha:chapter5} as necessary to achieve the tractability of the corridor mutex protocol.

    \item \textbf{A peer-to-peer spatial mutex protocol for narrow-aisle mutual exclusion.} Corridor contention is regulated by a lease-issuing protocol with no governing authority, using only short TTLs and best-effort broadcasting to renew a robot's right to occupy a narrow aisle. The absence of a coordination layer introduces opportunities for inconsistency, but the simplicity of a first-come-first-served arbitration policy makes the cost of potential coverage loss acceptable to the constraints of Chapter~\ref{cha:chapter5}, which analyzes the coverage loss of a similar protocol in detail. Chapter~\ref{cha:chapter8} evaluates the performance characteristics and deadlock-avoidance guarantees of the same protocol as a part of a broader eighteen-run campaign.

    \item \textbf{A hierarchical, two-stage recovery framework spanning localization, kinematic, and spatial-occlusion faults.} A single diagnostic watchdog covering an anisotropic covariance threshold triggers a spin-down and full-particle relocalization procedure for robot localization drift. A separate kinematic watchdog diagnoses stalls by querying position deltas rather than motor current, bypassing the lack of actuator feedback available to conventional low-level planners on an low-friction surface. Finally, a Pythagorean parallel-axis transformation lets a robot route around a completely failed, stationary peer without updating its allocation, preserving task continuity in the face of failures. Each is described in detail in Chapter~\ref{cha:chapter6}.

    \item \textbf{A battery-proportional, decentralized workload reallocation mechanism.} Every robot which observes a failed peer---through either of two diagnostic channels---updates its allocation to reflect a partitioning of the failed robot's targets proportional to its own battery level relative to all observing peers. Chapter~\ref{cha:chapter6} formally defines the update procedure, while Chapter~\ref{cha:chapter8} evaluates how much allocation-time overhead the procedure imposes on each robot in practice.

    \item \textbf{An edge-deployable vision-based anomaly classifier with closed-form explanation and analytical redundancy.} An explicit feature engineering pipeline of LBPs and HOGs fed into a linear SVM keeps each frame's inference within the processing budget of the robot's onboard planner while still retaining diagnostic value; a separate, structurally dissimilar XGBoost classifier is trained to secondary readiness in the event of a primary model failure. The same linearization of the vision stack which improves inference speed provides a secondary benefit of speeding up the attribution process by allowing gradient-based heatmapping to be replaced by a matrix projection. The pipeline is discussed in Chapter~\ref{cha:chapter5}, with analytical redundancy explored further in Chapter~\ref{cha:chapter6}.

    \item \textbf{A gate-based deterministic launch architecture for multi-robot ROS~2 bring-up.} Constrained by a monolithic launch file and without access to a central planner, the fleet brings up each component after the dependencies using four readiness-waiting primitives, blocking each stage of the bring-up process until the Gazebo world has loaded, the clock publisher has announced the current simulation time, the expected topics have appeared on the parameter server, and the necessary lifecycle nodes have transitioned to the active state. This replaces an error-prone series of hardcoded sleep statements from an earlier version of the launch file, as discussed in Chapter~\ref{cha:chapter7}.

    \item \textbf{An eighteen-run empirical campaign evaluating throughput, failure-to-reallocation latency, and coverage completeness as a function of fleet size.} Chapter~\ref{cha:chapter8} isolates the effect of fleet size from other variables which affect simulation performance in a nonlinear way, measures explicit reallocation latencies for a number of failure injection timelines, and identifies the cause of a coverage loss from an earlier campaign, described in contribution 2, to a diagnostic limitation in the blacklist propagation protocol.
\end{enumerate}

This is not a list of features which might be optionally deployed in a central planner; each of these contributions serves as a solution to a problem introduced by the specific constraints of a masterless fleet. As such, each is individually justified in later chapters.

\section{Thesis Organization\label{sec:ch1_organization}}

The remainder of this thesis is organized in two broad threads, separated roughly by Chapter~\ref{cha:chapter4}: the former provides the theoretical and empirical foundation the design choices are made upon, and the latter develops and analyzes the system itself.

Chapter~\ref{cha:chapter2} gathers the theoretical machinery used throughout the paper, including recursive Bayesian estimation and its particle filter incarnation, the non-holonomic dynamics of a wheeled robot, layered costmap-based planning and tracking control, distributed mutual exclusion algorithms and deadlock theory, fault classification, multi-robot task allocation taxonomy, and the statistical pattern recognition and explainable AI foundations of the vision pipeline used in the fleet autonomy stack, outside of the warehouse domain. Chapter~\ref{cha:chapter3} synthesizes the relevant literature on multi-robot task allocation, decentralized spatial coordination, fault-tolerant fleets, AMCL-based localization, and edge vision-based inspection systems, providing the context to the design decisions of this thesis, rather than contrasting it against a theoretical optimum.

Chapter~\ref{cha:chapter4} describes the coverage problem, mechanical limitations, and the simulation environment the remainder of the thesis is built upon. Chapters~\ref{cha:chapter5} and~\ref{cha:chapter6} develop the two main aspects of the system: the former focuses on the nominal operations, robot energy consumption, spatial coordination, the SDF-based waypoint geometry, path following controller, and the perception stack, while the latter details the system responses to faults, localization failures, path tracking slippage, vision-based replanning around occlusions, blacklist propagation, and battery-aware task redistribution. Chapter~\ref{cha:chapter7} focuses on the implementation specifics, package structure, deterministic launch files, TF-tree namespace isolation, diagnostics, and the operator interface, necessary to execute the algorithms described in Chapters~\ref{cha:chapter5} and~\ref{cha:chapter6}. Chapter~\ref{cha:chapter8} analyzes the results of an eighteen-run empirical evaluation, while Chapter~\ref{cha:chapter9} concludes the thesis, outlining the lessons learned during the development and highlighting the failure modes the implementation has been tested against, as well as the practical implications of the findings.